\subsection{Implementação dos Contadores}

Nesta etapa, as funções de transição mais básicas são concretizadas em contadores síncronos modulares: $mod6$, $mod10$ e $mod24$.

Cada contador é definido por uma tabela verdade derivada diretamente de sua transição de estados, sendo posteriormente simplificada por mapas de Karnaugh.

O sinal \textit{Carry} é empregado para indicar a ocorrência de overflow (no incremento) ou underflow (no decremento), dependendo da operação considerada.

A partir desse procedimento, a aplicação aos diferentes módulos segue de forma sistemática, variando apenas os domínios de cada contador.

\subsection{Contador Síncrono MOD $6$}

Seja $D$ o estado atual do contador, restrito ao conjunto de valores válidos.

\paragraph{Incremento}
A operação de incremento é definida por:
\[
D \mapsto (D + 1) \bmod 6
\]

A Tabela~\ref{tab:inc_mod6} apresenta a função de transição.

\begin{table} [H] \label{tab:dec_mod6}
\begin{center}
\begin{tabular}{c|ccc}
    $D$ & $O$ & Carry \\
    \hline
    $000$ & $001$ & $0$ \\
    $001$ & $010$ & $0$ \\
    $010$ & $011$ & $0$ \\
    $011$ & $100$ & $0$ \\
    $100$ & $101$ & $0$ \\
    $101$ & $000$ & $1$ \\
    $110$ & $000$ & $0$ \\
    $111$ & $000$ & $0$ \\    
\end{tabular}
\end{center}
\end{table}
    

As expressões booleanas minimizadas são:

\begin{align*}
    O_2 &= (\lnot D_2 \land D_1 \land D_0) \lor (D_2 \land \lnot D_1 \land \lnot D_0) \\[4pt]
    O_1 &= (\lnot D_2 \land \lnot D_1 \land D_0) \lor (\lnot D_2 \land D_1 \land \lnot D_0) \\[4pt]
    O_0 &= (\lnot D_2 \land \lnot D_0) \lor (\lnot D_1 \land \lnot D_0) \\[4pt]
    Carry &= D_2 \land \lnot D_1 \land D_0
\end{align*}


O sinal \textit{Carry} indica overflow.

As Figuras~\ref{fig:incmod6.png} e~\ref{fig:incmod6-red.png} apresentam a implementação do circuito.

\imgbig{incmod6.png}{Circuito lógico para incremento em módulo 6.}
\imgbig{incmod6-red.png}{Implementação em redstone do circuito de incremento em módulo 6.}

\paragraph{Decremento}
A operação de decremento é definida por:
\[
D \mapsto (D - 1) \bmod 6
\]

A Tabela~\ref{tab:dec_mod6} apresenta a função de transição.

\begin{table} [H] \label{tab:dec_mod6}
\begin{center}
\begin{tabular}{c|ccc}    
    $D$ & $O$ & Carry \\
    \hline
    $000$ & $101$ & $0$ \\
    $001$ & $000$ & $0$ \\
    $010$ & $001$ & $0$ \\
    $011$ & $010$ & $0$ \\
    $100$ & $011$ & $0$ \\
    $101$ & $000$ & $1$ \\
    $110$ & $000$ & $0$ \\
    $111$ & $000$ & $0$ \\
\end{tabular}
\end{center}
\end{table}

\begin{align*}
    O_2 &= (\lnot D_2 \land \lnot D_1 \land \lnot D_0) \lor (D_2 \land D_0) \\[4pt]
    O_1 &= (D_1 \land D_0) \lor (D_2 \land \lnot D_1 \land \lnot D_0) \\[4pt]
    O_0 &= \lnot D_0 \\[4pt]
    Carry &= \lnot D_2 \land \lnot D_1 \land \lnot D_0
\end{align*}

O sinal \textit{Carry} indica underflow.

As Figuras~\ref{fig:decmod6.png} e~\ref{fig:decmod6-red.png} mostram a implementação do circuito.

\imgbig{decmod6.png}{Circuito lógico para decremento em módulo 6.}
\imgbig{decmod6-red.png}{Implementação em redstone do circuito de decremento em módulo 6.}
\subsection{Contador Síncrono MOD $10$}

Seja $D$ o estado atual do contador, restrito ao conjunto de valores válidos.

\paragraph{Incremento}
A operação de incremento é definida por:
\[
D \mapsto (D + 1) \bmod 10
\]

A Tabela~\ref{tab:inc_mod10} apresenta a função de transição.

\begin{table} [H] \label{tab:dec_mod10}
\begin{center}
\begin{tabular}{c|cc}
    $D$ & $O$ & Carry \\
    \hline
    $0000$ & $0001$ & $0$ \\
    $0001$ & $0010$ & $0$ \\
    $0010$ & $0011$ & $0$ \\
    $0011$ & $0100$ & $0$ \\
    $0100$ & $0101$ & $0$ \\
    $0101$ & $0110$ & $0$ \\
    $0110$ & $0111$ & $0$ \\
    $0111$ & $1000$ & $0$ \\
    $1000$ & $1001$ & $0$ \\
    $1001$ & $0000$ & $1$ \\
    $\vdots$ & $\vdots$ & $\vdots$ \\
\end{tabular}
\end{center}
\end{table}

As expressões booleanas minimizadas são:

\begin{align*}
    O_3 &= (\lnot D_3 \land D_2 \land D_1 \land D_0) \lor (D_3 \land \lnot D_2 \land \lnot D_1 \land \lnot D_0) \\[4pt]
    O_2 &= (\lnot D_3 \land \lnot D_2 \land D_1 \land D_0) \lor (\lnot D_3 \land D_2 \land \lnot D_1) \lor (\lnot D_3 \land D_2 \land \lnot D_0) \\[4pt]
    O_1 &= (\lnot D_3 \land \lnot D_1 \land D_0) \lor (\lnot D_3 \land D_1 \land \lnot D_0) \\[4pt]
    O_0 &= (\lnot D_3 \land \lnot D_0) \lor (\lnot D_2 \land \lnot D_1 \land \lnot D_0) \\[4pt]
    Carry &= D_3 \land \lnot D_2 \land \lnot D_1 \land D_0
\end{align*}

O sinal \textit{Carry} indica overflow.

As Figuras~\ref{fig:incmod10.png} e~\ref{fig:incmod10-red.png} apresentam, respectivamente, a implementação lógica do circuito e sua realização em redstone.

\imgbig{incmod10.png}{Circuito lógico para incremento em módulo 10.}
\imgbig{incmod10-red.png}{Implementação em redstone do circuito de incremento em módulo 10.}

\paragraph{Decremento}

A operação de decremento é definida por:
\[
D \mapsto (D - 1) \bmod 10
\]

A Tabela~\ref{tab:dec_mod10} apresenta a função de transição.

\begin{table} [H] \label{tab:dec_mod10}
\begin{center}
\begin{tabular}{c|cc}
    $D$ & $O$ & Carry \\
    \hline
    $0000$ & $1001$ & $1$ \\
    $0001$ & $0000$ & $0$ \\
    $0010$ & $0001$ & $0$ \\
    $0011$ & $0010$ & $0$ \\
    $0100$ & $0011$ & $0$ \\
    $0101$ & $0100$ & $0$ \\
    $0110$ & $0101$ & $0$ \\
    $0111$ & $0110$ & $0$ \\
    $1000$ & $0111$ & $0$ \\
    $1001$ & $1000$ & $0$ \\
    $\vdots$ & $\vdots$ & $\vdots$ \\
\end{tabular}
\end{center}
\end{table}

\begin{align*}
    O_3 &= (\lnot D_3 \land \lnot D_2 \land \lnot D_1 \land \lnot D_0) \lor (D_3 \land D_0) \\[4pt]
    O_2 &= (D_2 \land D_0) \lor (D_2 \land D_1) \lor (D_3 \land \lnot D_0) \\[4pt]
    O_1 &= (D_1 \land D_0) \lor (D_2 \land \lnot D_1 \land \lnot D_0) \lor (D_3 \land \lnot D_0) \\[4pt]
    O_0 &= \lnot D_0 \\[4pt]
    Carry &= \lnot D_3 \land \lnot D_2 \land \lnot D_1 \land \lnot D_0
\end{align*}

O sinal \textit{Carry} indica underflow.

As Figuras~\ref{fig:decmod10.png} e~\ref{fig:decmod10-red.png} mostram a implementação lógica e sua versão em redstone.

\imgbig{decmod10-red.png}{Implementação em redstone do circuito de decremento em módulo 10.}
\imgbig{decmod10.png}{Circuito lógico para decremento em módulo 10.}

\subsection{Contador Síncrono MOD $24$}

Seja $D$ o estado atual do contador, restrito ao conjunto de valores válidos.

\paragraph{Incremento}
A operação de incremento é definida por:
\[
D \mapsto (D + 1) \bmod 24
\]

A Tabela~\ref{tab:inc_mod24} apresenta a função de transição.

\begin{table} [H] \label{tab:dec_mod24}
\begin{center}
\begin{tabular}{c|cc}
    $D$ & $O$ & Carry \\
    \hline
    $00\,0000$ & $00\,0001$ & $0$ \\
    $00\,0001$ & $00\,0010$ & $0$ \\
    $\vdots$ &  &  \\
    $00\,1001$ & $01\,0000$ & $0$ \\
    $01\,0000$ & $01\,0001$ & $0$ \\
    $\vdots$ &  &  \\
    $01\,1001$ & $10\,0000$ & $0$ \\
    $10\,0000$ & $10\,0001$ & $0$ \\
    $\vdots$ &  &  \\
    $10\,0010$ & $10\,0011$ & $0$ \\
    $10\,0011$ & $00\,0000$ & $1$ \\
\end{tabular}
\end{center}
\end{table}

As expressões booleanas minimizadas são:

\begin{align*}
    o_{11} &= (D_{10} \land D_{03} \land D_{00}) \lor (D_{11} \land \lnot D_{01}) \lor (D_{11} \land \lnot D_{00}) \\[4pt]
    o_{10} &= (\lnot D_{10} \land D_{03} \land D_{00}) \lor (D_{10} \land \lnot D_{03}) \lor (D_{10} \land \lnot D_{00}) \\[4pt]
    o_{03} &= (D_{02} \land D_{01} \land D_{00}) \lor (D_{03} \land \lnot D_{00}) \\[4pt]
    o_{02} &= (\lnot D_{11} \land \lnot D_{02} \land D_{01} \land D_{00}) \lor (D_{02} \land \lnot D_{01}) \lor (D_{02} \land \lnot D_{00}) \\[4pt]
    o_{01} &= (\lnot D_{03} \land \lnot D_{01} \land D_{00}) \lor (D_{01} \land \lnot D_{00}) \\[4pt]
    o_{00} &= \lnot D_{00} \\[4pt]
    Carry &= D_{11} \land D_{01} \land D_{00}
\end{align*}

O sinal \textit{Carry} indica overflow.

As Figuras~\ref{fig:incmod24} e~\ref{fig:incmod24-red} apresentam, respectivamente, a implementação lógica do circuito e sua realização em redstone.

\imgbig{incmod24.png}{Circuito lógico para incremento em módulo 24.}
\imgbig{incmod24-red.png}{Implementação em redstone do circuito de incremento em módulo 24.}

\paragraph{Decremento}
A operação de decremento é definida por:
\[
D \mapsto (D - 1) \bmod 24
\]

A Tabela~\ref{tab:dec_mod24} apresenta a função de transição.

\begin{table} [htbp] \label{tab:dec_mod24}
\begin{center}
\begin{tabular}{c|cc}
    $D$ & $O$ & Carry \\
    \hline
    $00\,0000$ & $10\,0011$ & $1$ \\
    $00\,0001$ & $00\,0000$ & $0$ \\
    $\vdots$ &  &  \\
    $00\,1001$ & $00\,1000$ & $0$ \\
    $01\,0000$ & $00\,1001$ & $0$ \\
    $\vdots$ &  &  \\
    $01\,1001$ & $01\,1000$ & $0$ \\
    $10\,0000$ & $01\,1001$ & $0$ \\
    $\vdots$ &  &  \\
    $10\,0010$ & $10\,0001$ & $0$ \\
    $10\,0011$ & $10\,0010$ & $0$ \\
\end{tabular}
\end{center}
\end{table}

As expressões booleanas minimizadas são:

\begin{align*}
    o_{11} &= (\lnot D_{11} \land \lnot D_{10} \land \lnot D_{03} \land \lnot D_{02} \land \lnot D_{01} \land \lnot D_{00}) \\
    &\quad \lor (D_{11} \land D_{00}) \lor (D_{11} \land D_{01}) \\[4pt]
    o_{10} &= (D_{10} \land D_{00}) \lor (D_{10} \land D_{01}) \lor (D_{10} \land D_{02}) \lor (D_{10} \land D_{03}) \\
    &\quad \lor (D_{11} \land \lnot D_{01} \land \lnot D_{00}) \\[4pt]
    o_{03} &= (D_{03} \land D_{00}) \lor (D_{10} \land \lnot D_{03} \land \lnot D_{02} \land \lnot D_{01} \land \lnot D_{00}) \\
    &\quad \lor (D_{11} \land \lnot D_{01} \land \lnot D_{00}) \\[4pt]
    o_{02} &= (D_{02} \land D_{00}) \lor (D_{02} \land D_{01}) \lor (D_{03} \land \lnot D_{00}) \\[4pt]
    o_{01} &= (\lnot D_{11} \land \lnot D_{10} \land \lnot D_{01} \land \lnot D_{00}) \\
    &\quad \lor (D_{01} \land D_{00}) \lor (D_{02} \land \lnot D_{01} \land \lnot D_{00}) \\
    &\quad \lor (D_{03} \land \lnot D_{00}) \\[4pt]
    o_{00} &= \lnot D_{00} \\[4pt]
    Carry &= \lnot D_{11} \land \lnot D_{10} \land \lnot D_{03} \land \lnot D_{02} \land \lnot D_{01} \land \lnot D_{00}
\end{align*}

O sinal \textit{Carry} indica underflow.

As Figuras~\ref{fig:decmod24} e~\ref{fig:decmod24-red} mostram a implementação lógica e sua versão em redstone.

\imgbig{decmod24.png}{Circuito lógico para decremento em módulo 24.}
\imgbig{decmod24-red.png}{Implementação em redstone do circuito de decremento em módulo 24.}

\subsubsection{Considerações de Projeto}

Note que, embora funcional, o contador apresentado não é autocorretivo, pois estados inválidos não são mapeados explicitamente para estados válidos. Isso implica que o sistema não garante convergência para o conjunto de estados válidos a partir de estados inválidos. Em aplicações reais, isso pode levar a comportamentos inesperados.

